\begingroup
\raggedbottom
\clubpenalty=10000
\widowpenalty=10000

\chapter{From Models to Systems}
\label{ch:models-systems}

A model produces an output; a system turns that output into consequences. A score may determine which cases receive attention, an answer may influence a person's decision, and a referral may join a queue. These choices change the effective behavior even when the fitted parameters remain fixed.

Write the served predictor as $f_\theta(\phi_v(x))$, where $\phi_v$ is a particular version of the input representation. An action is
\[
a=\delta\bigl(f_\theta(\phi_v(x)),c\bigr),
\]
where $c$ includes costs, available evidence, and resource constraints. Evaluating $f_\theta$ and evaluating this complete decision process are different problems. The second includes the measurements that reach the model, the actions that follow, and the outcomes that become observable.

Two familiar examples make the distinction concrete. A student-risk predictor estimates whether a student will miss at least one of two specified deadlines in weeks three and four, using information available at the end of week two. An advising team may have capacity to review only five students. A course-support assistant, by contrast, retrieves institutional policy and drafts answers to incoming questions; its referrals compete for advisers' time. Predictive error matters in both cases, but it does not describe the complete workflow.

\section{The Objective Belongs to the Decision Process}

Suppose the assistant continues to produce the same answers it produced during evaluation, but its retrieved policy documents become outdated and its referral queue doubles. The weights are unchanged; the evaluated composition is not. A useful requirement therefore names an outcome, a population, a measurement window, and a decision rule.

For example, a response-time target needs to specify whether time ends at the first automated acknowledgement or at a usable answer. A severe-error rate needs a denominator and a way of finding errors, including cases that do not attract complaints. A backlog target needs to distinguish the mean queue from the longest wait. Writing these requirements makes the intended behavior assessable; it does not establish that the service satisfies them.

Means and tail behavior answer different questions. If $80\%$ of requests finish in $5$ milliseconds and $20\%$ in $50$, the mean time is
\[
0.8(5)+0.2(50)=14\text{ milliseconds}.
\]
Using the smallest time by which at least $95\%$ of requests finish, the $95$th percentile is $50$ milliseconds. A satisfactory mean can coexist with a slow experience for many users.

\begin{figure}[t]
\centering
\begin{tikzpicture}[ml diagram, x=1cm, y=1cm,
  stage/.style={ml box, minimum width=1.8cm, minimum height=0.9cm}]
\node[stage] (input) at (0,0) {Measurement};
\node[stage] (predict) at (2.65,0) {Model};
\node[stage, ml key box] (policy) at (5.3,0) {Decision};
\node[stage] (action) at (7.95,0) {Consequence};
\draw[ml arrow] (input) -- (predict);
\draw[ml arrow] (predict) -- (policy);
\draw[ml arrow] (policy) -- (action);
\node[ml label] (context) at (5.3,1.05) {Costs and capacity};
\draw[ml arrow] (context) -- (policy);
\draw[ml feedback] (action.south) -- (7.95,-1.15) -- (0,-1.15) -- (input.south);
\node[ml label, fill=white, inner sep=3pt] at (3.975,-1.15) {later observations};
\end{tikzpicture}
\caption{The decision process includes representation, action, and feedback. Later observations may depend on earlier actions.}
\label{fig:system-pipeline}
\end{figure}

\subsection{Costs and Capacity Change the Preferred Action}

Consider an urgent-case classifier with $p_i=P(Y_i=1\mid x_i)$. Assume that incorrectly flagging a non-urgent case costs $1$, failing to flag an urgent case costs $9$, and correct decisions cost zero. These are stipulated decision costs, not estimates of the causal benefit of contacting a student. Flagging has conditional expected loss $1-p_i$; not flagging has loss $9p_i$. Flagging is preferred when
\[
1-p_i<9p_i,\qquad\text{or}\qquad p_i>0.1.
\]
At equality, either action has the same expected loss.

Now let $u_i\in\{0,1\}$ indicate whether case $i$ is flagged, with at most $B$ flags:
\[
\min_{u_1,\ldots,u_n}\ \sum_i\left[u_i(1-p_i)+(1-u_i)9p_i\right],
\qquad \sum_i u_i\le B.
\]
Relative to flagging nobody, the expected loss reduction for case $i$ is $10p_i-1$. With equal resource use, choose the largest positive reductions until capacity is exhausted. Linearity of expectation makes this calculation valid even if the labels are dependent. It still assumes additive costs and no effects of one allocation on another case's loss.

For probabilities $(0.8,0.3,0.12,0.08)$ and capacity $B=2$, the reductions are $(7,2,0.2,-0.2)$. Flagging nobody has expected loss $9(1.3)=11.7$. Flagging the first two cases lowers it to $2.7$. A capacity limit changes the operating rule even though every probability is unchanged.

Probability alone does not rank benefits when costs differ. With false-negative cost $20$, false-positive cost $1$, and $p=0.2$, the reduction is $20(0.2)-0.8=3.2$. With costs $9$ and $1$ and $p=0.3$, it is $2.7-0.7=2$. The lower-probability case has greater expected benefit under those costs.

If cases consume different amounts of time, the constraint becomes $\sum_i d_i u_i\le B$, where $d_i$ is the resource required. Selecting indivisible cases is then a knapsack problem; ordering by benefit per minute need not yield the optimum. A mandatory human-review requirement also cannot be relaxed merely to make a queue fit. The feasible response may be to restrict the service's scope, add capacity, or use an authorized alternative.

\section{Referrals Create a Queue}
\label{subsec:queue-capacity-threshold}

Let $q_t$ be the backlog at the start of day $t$, $a_t$ the new cases arriving before that day's service, and $r_t$ the number that can be completed that day. In a toy model with equal-sized cases, no abandonment, and service whenever work is available,
\[
q_{t+1}=\max\{0,q_t+a_t-r_t\}.
\]
The order of arrivals and service is part of the model. It does not describe within-day waiting times.

With $q_0=0$, capacity $10$, and arrivals $12,8,14,6,10$, the successive backlogs are $2,0,4,0,0$. Mean arrivals equal capacity, yet a queue forms on some days. With capacity $30$ and exactly $36$ arrivals every day, $q_t=6t$; with exactly $24$, the day-end backlog stays zero. Replacing fluctuating arrivals by their mean would turn these deterministic illustrations into an unjustified forecast.

In particular, expectation does not commute with the positive-part operation. If the first day's arrivals are $0$ or $20$ with equal probability and capacity is $10$, then
\[
E[q_1]=\tfrac12(0)+\tfrac12(10)=5,
\qquad
\max\{0,E[a_0]-10\}=0.
\]
An average load below average capacity can support stability under suitable assumptions, but does not imply an empty queue, a bound on every delay, or even finite mean waiting time for every arrival process. Equal average load and capacity are especially insufficient as a stability guarantee.

For a stable process with finite, consistently defined long-run averages, Little's law states
\[
L=\lambda W,
\]
where $L$ is average occupancy, $\lambda$ is the effective arrival rate, and $W$ is average time spent in that part of the process \citep{little1961}. Queue occupancy goes with waiting time; total occupancy goes with total time, including service. If an average of $12$ cases wait and $24$ cases enter the queue per day, average waiting time is $12/24=0.5$ days. This identity neither establishes stability nor bounds an individual case's delay.

Capacity affects statistical policy through these consequences. Lowering a referral threshold may increase recognition of difficult cases while increasing delays for everyone. Choosing a threshold requires evaluating both effects under a defined fallback policy. Automatically answering cases that require review is not justified by a calculation showing that the alternative queue is too large.

\section{The Served Representation Must Preserve Meaning}

Training-serving skew is a mismatch between the representation evaluated during development and the representation used to make decisions. Units, missing-value meanings, time cutoffs, and available evidence are all part of that representation.

If training uses $z=(x-\mu_{\rm tr})/\sigma_{\rm tr}$ and serving uses $\tilde z=(x-\mu_{\rm sv})/\sigma_{\rm sv}$, with positive scales, then
\[
\tilde z=
\frac{\sigma_{\rm tr}}{\sigma_{\rm sv}}z+
\frac{\mu_{\rm tr}-\mu_{\rm sv}}{\sigma_{\rm sv}}.
\]
For the score $s=z-0.5$, training constants $(\mu_{\rm tr},\sigma_{\rm tr})=(50,10)$ put the raw-input boundary at $55$. Serving constants $(51,11)$ move it to $56.5$. At $x=56$, the evaluated score is $0.1$ but the served score is approximately $-0.0455$. Fixed weights can therefore give a different decision after a preprocessing change.

A time cutoff is just as important as a numerical scale. A week-two predictor cannot use an adviser note written in week four, even if that note is present in the historical database. A missing attendance measurement is not automatically zero attendance. Replacing missing values by zero changes the meaning unless the model and its evaluation explicitly support that convention.

For a retrieval-based assistant, the evidence contract includes the relevant institution, course, term, effective date, and access permissions. An unchanged language model can answer differently when the index contains an obsolete policy. A versioned record of the model, preprocessing, decision rule, and evidence collection helps identify what was actually evaluated. Restoring an earlier combination can restore its former behavior; an old policy may nevertheless be inappropriate after the institution changes its rules.

Sculley et al.\ analyze dependencies and feedback that accumulate around learned models \citep{sculley2015}. Amershi et al.'s organizational case study describes the importance of data management and validation in machine-learning workflows \citep{amershi2019}. These sources explain mechanisms by which a successful model can become part of an unreliable service. They do not establish that representation errors dominate every deployment.

\begin{figure}[t]
\centering
\begin{tikzpicture}[ml diagram, x=1cm, y=1cm,
  stage/.style={ml box, text width=2.45cm, minimum height=1.1cm}]
\node[stage] (evidence) at (0,0) {Eligible evidence\\and measurements};
\node[stage, ml key box] (predictor) at (3.55,0) {Served predictor\\and representation};
\node[stage] (decision) at (7.1,0) {Answer or refer\\within capacity};
\draw[ml arrow] (evidence) -- (predictor);
\draw[ml arrow] (predictor) -- (decision);
\node[stage] (outcomes) at (7.1,-2.0) {Observed outcomes\\and waiting time};
\node[stage] (change) at (0,-2.0) {Reviewed change\\and new evaluation};
\draw[ml arrow] (decision) -- (outcomes);
\draw[ml arrow] (outcomes) -- (change);
\draw[ml feedback] (change) -- (evidence);
\node[ml label, fill=white, inner sep=3pt] at (3.55,-2.0) {review evidence};
\end{tikzpicture}

\caption{The evaluated system is a compatible combination of evidence, predictor, and decision policy. Reviewed outcomes can motivate a change, whose effects must then be evaluated.}
\label{fig:complete-system-architecture}
\end{figure}

\section{Evaluate the Model and Its Fallback Together}

A selective predictor answers some cases and defers others. Let $A(x)\in\{0,1\}$ indicate an automated answer, $\ell_m$ its loss, and $\ell_h$ the loss from the fallback process. Coverage and selective risk are
\[
\phi=E[A(X)],\qquad
R_{\rm sel}=\frac{E[A(X)\ell_m]}{\phi},\quad \phi>0.
\]
At zero coverage this ratio is undefined, not evidence of zero error. With referral cost $c$ per deferred case, the complete expected loss is
\[
R_{\rm sys}
=E\!\left[A(X)\ell_m+(1-A(X))\ell_h\right]
+cE[1-A(X)].
\]
The fallback's performance must be measured on the cases actually sent to it, under the interface and workload it actually encounters. A human's accuracy on an unrelated collection of easy cases does not supply that quantity.

Suppose $60$ of $100$ cases are answered automatically, with $3$ errors; the other $40$ receive human answers, with $4$ errors. Coverage is $0.6$, selective error is $3/60=0.05$, and complete error is $7/100=0.07$. If each referral costs $0.02$ in the same loss units, total empirical loss is $0.07+0.02(0.4)=0.078$.

To compare with answering all cases automatically, we also need the candidate automated answers on the $40$ deferred cases. If a valid audit finds $12$ errors there, full automation would have $15$ errors out of $100$. The $5\%$ error on the accepted cases alone cannot establish this counterfactual comparison. Evaluating selective classification explicitly separates coverage from conditional risk \citep{geifman2017}.

\subsection{Defer Cases Where the Alternative Helps}

Let $r_m(x)$ and $r_h(x)$ be conditional expected losses for the model and fallback. Referring case $x$ reduces expected loss by
\[
b(x)=r_m(x)-r_h(x)-c.
\]
With equal review times and a fixed batch budget, the optimal additive allocation selects the largest positive $b(x)$ values. The least confident model prediction need not have the largest benefit from human review.

For example, suppose model and human error probabilities are $(0.40,0.35)$ for case A and $(0.20,0.02)$ for case B. With zero review cost, the gains are $0.05$ and $0.18$. Reviewing B helps more even though the model is less likely to be wrong there. With cost $0.1$, only B has positive gain; with cost $0.2$, neither does. These conclusions rely on knowing the relevant conditional error rates and on the additive decision model.

For a fixed score, raising an acceptance threshold reduces coverage, apart from tie conventions. It does not guarantee that conditional error falls at every threshold: the score may order errors imperfectly or behave differently across subgroups. Calibration of a probability score also does not establish that an individual accepted case is safe.

\section{Human Review Is a Conditional Decision}

Displaying a draft changes the task performed by the reviewer. A person may accept a plausible mistake, spend time checking a correct answer, or overlook missing evidence. Review therefore belongs in the evaluated intervention, rather than entering the loss calculation as a perfect final step.

A useful comparison measures independent outcome labels, review time, and referral consequences under the actual interfaces. Collecting an initial judgment before displaying the model's answer can reveal anchoring effects, but this interface choice does not by itself guarantee independent or correct judgments. The comparison still needs an appropriate assignment and measurement design.

Probability displays should name the event and relevant population. A statement about missing at least one of the two specified deadlines during weeks three and four means something different from a statement about failing the course. A value such as $70\%$ also needs a calibration context; it is not certainty about the person currently being reviewed. Lipkus's review of numerical, verbal, and visual risk communication describes evidence and unresolved questions, rather than a single universally superior format \citep{lipkus2007}.

An attribution can help a reviewer inspect what affected a score, but it is not automatically a causal explanation or a justification for the resulting action. Likewise, an override rate is not an accuracy measure. If reviewers accept $19$ of $20$ recommendations, acceptance is $95\%$; without independent labels, their accuracy remains unknown. A low or high override rate has no universal optimum.

The downstream process can be physical rather than organizational. In a pedestrian-alert system, response time includes sensing, prediction, control, actuation, and braking. Faster inference alone does not establish that the complete stopping process meets its requirement. The relevant delay budget must follow the physical task and operating conditions.

\section{Monitoring Measures Change Before It Explains It}

Monitoring can distinguish input validity, changes in predictions, referral load, and measured outcomes. These observations answer different questions. A change in an input distribution does not establish an increase in error, and stable inputs do not establish stable performance.

\subsection{Marginal Comparisons Have Limited Scope}

For two distributions over the same bins, with strictly positive proportions $p_b$ and $q_b$, the population stability index is
\[
\operatorname{PSI}(p,q)=
\sum_b(q_b-p_b)\log\frac{q_b}{p_b}.
\]
This is the symmetric sum $D_{\rm KL}(q\|p)+D_{\rm KL}(p\|q)$ for those binned distributions. The identity does not make a particular alert threshold universal. Bin boundaries, sample sizes, and smoothing of zero counts affect the measured value. For $p=(0.6,0.4)$ and $q=(0.3,0.7)$,
\[
\operatorname{PSI}=0.3\log(3.5)\approx0.3758.
\]
The value describes a distributional discrepancy, not an error rate.

The two-sample Kolmogorov--Smirnov statistic compares empirical cumulative distributions:
\[
D_{n,m}=\sup_z\left|\widehat F_n(z)-\widehat G_m(z)\right|.
\]
For samples $(0,0,1,1)$ and $(0,1,1,2)$, it is $0.25$. Ties do not make this statistic undefined, but ordinary continuous-null significance calculations do not automatically apply unchanged to discrete data. Dependence among observations and repeated monitoring also matter when interpreting significance \citep{nistKS}.

Marginal monitoring can miss a joint change completely. Let $X_1,X_2$ be binary. Under distribution $P$, give probability $1/2$ to each of $(0,0)$ and $(1,1)$; under $Q$, give probability $1/2$ to each of $(0,1)$ and $(1,0)$. Both feature marginals are identical. If the target is $Y=X_1\mathbin{\mathrm{XOR}}X_2$ and the predictor always returns zero, its error rises from $0$ under $P$ to $1$ under $Q$. The prediction distribution is also unchanged. The example shows why particular monitoring signals cannot certify performance.

\subsection{Labels Need an Exposure and Observation Window}

A future-deadline label is not complete before the stated future period ends. Evaluating newly scored students before their deadlines occur can misclassify unresolved outcomes. Complaints about an assistant are also a selected signal: silence does not establish correctness.

The decision process can determine which labels become visible. If only referred cases receive careful review, the observed error rate concerns that selected population. Chapter~\ref{ch:reliability}'s distinction between complete and conditional risk applies here. Chapter~\ref{ch:ranking-recommendation} develops exposure and logged-policy identification assumptions. A planned audit can broaden observation, but its sampling probabilities, target population, and outcome window must be specified.

An alert prompts investigation; it does not identify the remedy. A broken unit conversion calls for repairing the representation. A corrupted label collection calls for repairing measurement. Retraining on either failure can preserve the error. Restoring an earlier model is useful only if its representation, evidence, and decision policy are still compatible with the current task.

\subsection{Security Sets the Allowed Actions}
\label{sec:security-boundaries}

The security distinctions in Chapter~\ref{ch:reliability} continue to apply to the complete service. Text from a retrieved document can contain instructions, but its presence in context does not authorize an action. The actions permitted to the assistant need an authority boundary independent of what the text requests.

Limits on request rate can reduce some repeated-query attacks, but do not establish protection against a single revealing response or a distributed attacker. Size and schema checks do not establish that retrieved content is trustworthy. A distributional alert can accompany an attack, an ordinary policy change, or a measurement error; the absence of an alert does not establish security.

Response options depend on the affected boundary: restrict a capability, remove compromised evidence, restore a compatible combination, or route affected cases to an authorized fallback. Records needed for investigation also need a defined access and retention policy. Logging everything indefinitely is not a mathematical or operational guarantee of reliability.

\section{A Rollout Is Also an Experiment}

In shadow evaluation, a candidate receives live inputs while another policy continues to determine actions. This can reveal incompatible representations, computational costs, and disagreement. It cannot generally reveal outcomes under the candidate's actions, because those actions did not occur.

A limited rollout exposes fewer cases to a new policy. This can reduce the number affected by a failure, but does not bound its severity or show that rare failures are absent. A selected pilot group is not automatically representative or randomized.

\subsection{Randomization Must Match the Outcome Unit}

For a randomized comparison, define the assigned alternatives, the unit of assignment, the target population, and the observation window. The difference in mean outcomes estimates the effect of assignment under the design's assumptions. It need not estimate the effect of actual use when people ignore or switch their assigned alternatives.

Shared capacity creates possible interference. Assigning individual students to different referral policies while all referrals enter one adviser queue can make one student's assignment affect another student's delay. Cluster assignment or time blocks may help align the design with that mechanism, but introduce their own assumptions about within-cluster dependence, time trends, and carryover.

Power calculations depend on the same design. For an approximate two-sided comparison of independent Bernoulli outcomes near probability $0.5$, equal group sizes, $5\%$ significance, and $80\%$ power, detecting a difference of $0.01$ gives
\[
n\approx
\frac{2(1.96+0.84)^2(0.5)^2}{(0.01)^2}
=39{,}200
\]
observations per group. Here $0.5$ is the outcome standard deviation, not the standard error. Clustering, missing outcomes, and other departures change the calculation, as discussed in Chapter~\ref{ch:experiments-claims}.

Repeatedly checking a fixed-horizon test and stopping at the first favorable result does not preserve its nominal false-positive rate. A prespecified sequential analysis can accommodate interim decisions; an operational stop for a severe failure is a different decision from a claim of statistical superiority. Observing no severe failures in a finite sample also does not establish a zero population rate.

\subsection{Adaptive Policies Change What Is Observed}

A bandit policy changes its action probabilities as feedback arrives. Its records need to connect contexts, available actions, chosen actions, probabilities, and outcomes. Delayed labels and changing eligibility affect the interpretation.

Exploration can improve coverage of alternatives, but an exploration bonus does not guarantee that every relevant action is sampled sufficiently. Existing evidence may already support a comparison; exploring solely for its own sake is not a universal requirement. When action support and identification assumptions hold, weighting can identify a fixed policy's expected outcome from suitable logged feedback. Partial observation does not by itself make identification approximate. Finite-sample uncertainty, unsupported actions, and changing behavior remain separate concerns.

\section{Saving Resources Creates a New Comparison}
\label{sec:production-infra}

A shared feature definition or a registry of evaluated combinations can help keep a system coherent. These are mechanisms for preserving meaning and tracking evidence; no particular infrastructure is required by the mathematics. A simpler or smaller model is also a new candidate whose full behavior needs comparison.

\subsection{Distillation Transfers a Distribution}

Knowledge distillation trains a student using a teacher's outputs \citep{hinton2015distillation}. For teacher and student logits $z^T,z^S$ and temperature $T>0$, define
\[
q_j^T=\frac{\exp(z_j^T/T)}{\sum_k\exp(z_k^T/T)},
\qquad
q_j^S=\frac{\exp(z_j^S/T)}{\sum_k\exp(z_k^S/T)}.
\]
One objective, with $\alpha\in[0,1]$, is
\[
\mathcal L=
\alpha T^2D_{\rm KL}(q^T\|q^S)
+(1-\alpha)\operatorname{CE}(y,q^S_{\!1}),
\]
where $q^S_{\!1}$ uses temperature one. The teacher distribution is fixed during this student update. Replacing the KL term by teacher-to-student cross-entropy changes the objective only by a constant teacher entropy.

For the softened term alone,
\[
\frac{\partial\,[T^2D_{\rm KL}(q^T\|q^S)]}{\partial z_j^S}
=T(q_j^S-q_j^T).
\]
At large $T$, with finite logits and centered logits $\bar z_j=z_j-K^{-1}\sum_k z_k$, this approaches $(\bar z_j^S-\bar z_j^T)/K$. The $T^2$ factor compensates for an approximate high-temperature gradient scaling; it does not make gradients exactly temperature-independent.

For teacher logits $(\log4,0)$, the temperature-one distribution is $(0.8,0.2)$. At $T=2$, it is $(2/3,1/3)$. A student with equal logits has distribution $(1/2,1/2)$, and the gradient of the rescaled softened term is $(-1/3,1/3)$. Distillation encourages agreement with the teacher. It does not establish that the teacher is correct, that the student is equally accurate, or that either fits the service's time budget.

\subsection{Quantization and Pruning Need an Error Argument}

Rounding a scalar weight to the nearest multiple of step $s$ gives error at most $s/2$, provided no clipping occurs. For a linear score with exact input and bias,
\[
\left|(\tilde w-w)^\top x\right|
\le\frac{s}{2}\|x\|_1.
\]
If the original absolute score exceeds this bound, rounding cannot change its sign. The argument does not cover clipping, quantized inputs, or the accumulated effects of a nonlinear network without additional analysis.

Take $w=(0.26,-0.44)$, $x=(2,-1)$, and $s=0.1$. Rounding gives $\tilde w=(0.3,-0.4)$. The original score is $0.96$, the rounded score is $1.00$, and the actual change $0.04$ is below the bound $0.15$. Changing weights from $32$ to $8$ bits reduces their raw bit count to one quarter; auxiliary scales, storage layout, and other system costs are additional.

Pruning removes parameters or structures. A numerically small weight need not have a small effect: a weight $0.01$ multiplied by an input $1000$ contributes $10$. Structured removal can change the amount of computation a device performs. Unstructured sparsity saves time only when the representation and execution method can use it. Parameter count alone does not determine latency.

\subsection{A Cascade Trades Average Cost Against Conditional Behavior}

A cascade sends easy cases to a small model and routes selected cases to a larger one. Suppose every request first takes $5$ milliseconds, and routed cases take another $45$. At routing fraction $0.2$, mean time is $14$ milliseconds and the $95$th percentile is $50$. At routing fraction $0.02$, mean time is $5.9$ milliseconds and the $95$th percentile is $5$, under this exact two-point model.

The error calculation must also follow the route. Accuracy of each model on the whole population does not determine its accuracy on the cases selected for it. A faster cascade can fail if the small model confidently keeps precisely the cases on which it is wrong.

A cached answer has the same conceptual issue: reusing it assumes that the relevant context is unchanged. New policy evidence, altered access rights, or a different question can break that assumption. Resource savings should therefore be measured together with outcomes on the actual population, including the cases affected by the saving mechanism.

\section{Changes Need a Defined Scope and Owner}

A release specifies a task, time cutoff, population, evaluated combination, decision policy, and fallback. Someone needs authority to respond when these conditions fail. A defined owner makes a response possible; it does not establish that all failures will be detected.

Different findings call for different changes. An outdated source can require an evidence update. A missing measurement can require a revised representation and evaluation. An overloaded referral queue can require a narrower service or more capacity. Increasing model complexity does not automatically repair any of these conditions.

For student support, estimating future deadlines and deciding whom to contact remain distinct questions. A good prediction plus a five-case capacity calculation does not establish that outreach benefits students; that claim requires evidence about the intervention. For an assistant, a credible answer on a routine policy question does not establish readiness for every advising question. The defensible scope follows the evidence about the complete process.

\section{Study Questions}

\begin{enumerate}
\item Why can a fixed model file produce a different served decision after a representation change?
\item How do conditional prediction error, decision cost, and causal intervention benefit differ?
\item Why does a capacity limit favor expected loss reductions rather than probabilities alone?
\item Which assumptions make the daily queue recurrence appropriate?
\item What does Little's law establish, and which delay claims remain unsupported?
\item Why is selective error insufficient for evaluating the complete fallback system?
\item When would reviewing a relatively confident prediction help more than reviewing an uncertain one?
\item What can an override rate reveal without independent outcome labels?
\item How can feature and prediction marginals stay unchanged while error increases?
\item Which outcomes can shadow evaluation observe, and which require actions to occur?
\item Why does a shared queue complicate individual randomization?
\item What do distillation, quantization, and a cascade each change about the evaluated predictor?
\end{enumerate}

\section{Exercises}

\begin{enumerate}
\item \textbf{Allocate a fixed budget.} Use the four probabilities $(0.8,0.3,0.12,0.08)$, false-positive cost $1$, false-negative cost $9$, and capacity two. Assume additive costs and no interference. Derive each flag's expected loss reduction, select the cases, and calculate expected loss before and after allocation.\par

\item \textbf{Vary cost and duration.} Compare a case with $p=0.2$ and costs $(C_{\rm FP},C_{\rm FN})=(1,20)$ with a case with $p=0.3$ and costs $(1,9)$. Then consider three indivisible cases with benefits $(6,10,12)$ and durations $(1,2,3)$ under a duration budget of $5$. Show that selecting by benefit per unit time can miss the best combination.\par

\item \textbf{Follow a queue.} Starting with zero backlog, use capacity $10$ and arrivals $12,8,14,6,10$, arriving before service. Calculate each next-day backlog. Separately, let first-day arrivals be $0$ or $20$ with equal probability. Compare expected backlog with the backlog computed from expected arrivals.\par

\item \textbf{Keep time definitions consistent.} A stable service handles $24$ cases per day, with average queue occupancy $12$ and average total occupancy $18$. Use Little's law to calculate mean waiting time and mean total time. Explain why neither quantity is a maximum delay.\par

\item \textbf{Move a boundary.} A score is $z-0.5$. Training uses $z=(x-50)/10$; serving uses $(x-51)/11$. Find both raw-input boundaries and both scores at $x=56$.\par

\item \textbf{Account for the fallback.} Of $100$ cases, the system automates $60$ with $3$ errors and refers $40$ with $4$ human errors. Calculate coverage, selective error, and complete error. Add cost $0.02$ per referral. An independent audit finds $12$ errors in the candidate automated answers for the deferred cases; calculate full-automation error. State the additional evidence required to evaluate an all-human policy.\par

\item \textbf{Choose a review.} Model and fallback error probabilities are $(0.40,0.35)$ for A and $(0.20,0.02)$ for B. At most one case can be reviewed. Calculate the optimal choice for review costs $0$, $0.1$, and $0.2$, assuming unit error loss, equal review times, and known conditional rates.\par

\item \textbf{Measure a discrepancy.} Calculate PSI for $(0.6,0.4)$ and $(0.3,0.7)$. Calculate the empirical KS statistic for $(0,0,1,1)$ and $(0,1,1,2)$. Explain why neither result determines prediction error or supplies a universal alarm threshold.\par

\newpage
\item \textbf{Construct an invisible joint change.} Use the two binary-feature distributions in the monitoring example. Write each marginal, then calculate the error of the constant-zero predictor for target $X_1\mathbin{\mathrm{XOR}}X_2$. Which monitored distributions remain unchanged?\par

\item \textbf{Differentiate distillation.} Use teacher logits $(\log4,0)$, student logits $(0,0)$, temperature $2$, class-one label $y=(1,0)$, and $\alpha=1/2$. Calculate both softened distributions, $T^2D_{\rm KL}(q^T\|q^S)$, the complete objective, and its gradient with respect to the student logits. Use natural logarithms.\par

\item \textbf{Bound rounding error.} Round $(0.26,-0.44)$ to step $0.1$ and evaluate the score at $(2,-1)$, with zero bias. Compare the actual change with the $\ell_1$ bound. Explain which score margins are sufficient to preserve a binary sign decision under the stated assumptions.\par

\item \textbf{Inspect a cascade's tail.} Every request costs $5$ milliseconds, with another $45$ for routed cases. Calculate mean and $95$th-percentile latency at routing fractions $0.2$ and $0.02$. Identify the conditional error rates needed to compare the cascades.\par

\item \textbf{Evaluate the review interface.} Design a comparison of independent human answering, answering after viewing a draft, and selective human fallback. Define assignment, independent outcome measurement, review time, and missing outcomes. Explain why acceptance of a draft is not the outcome label.\par

\item \textbf{Specify a bounded release.} For the course-support assistant, define eligible questions, evidence validity, referral capacity, observation windows, and a fallback if an evidence source is withdrawn. State one predictive claim and one causal claim, and identify the different evidence each requires.\par

\end{enumerate}

\subsection*{Selected Answers}

\begin{enumerate}
\item The gains are $7,2,0.2,-0.2$. Choose the first two; expected loss falls from $11.7$ to $2.7$.
\item The gains are $3.2$ and $2$. Duration ratios are $6,5,4$: the greedy choice takes the first two for benefit $16$, while the second and third fit the same budget and give $22$.
\item Backlogs are $2,0,4,0,0$. In the random first-day example, expected backlog is $5$; using mean arrivals gives zero.
\item Mean waiting time is $0.5$ days and mean total time $0.75$ days.
\item The boundaries are $55$ and $56.5$; scores at $56$ are $0.1$ and $-1/22\approx-0.0455$.
\item Coverage is $0.6$, selective error $0.05$, complete error $0.07$, and complete loss with referral cost $0.078$. Full automation has error $0.15$. The all-human comparison needs human outcomes on the $60$ currently automated cases under the proposed review conditions.
\item Review B at costs $0$ and $0.1$; review neither at cost $0.2$.
\item PSI is $0.3\log3.5\approx0.3758$ and KS is $0.25$.
\item Each feature is Bernoulli with probability $1/2$ under both distributions. Error is $0$ under $P$ and $1$ under $Q$; both feature marginals and the constant prediction distribution are unchanged.
\item The teacher and student softened distributions are $(2/3,1/3)$ and $(1/2,1/2)$. The rescaled KL term is approximately $0.22653$, the complete objective approximately $0.45984$, and its gradient $(-5/12,5/12)$.
\item The scores are $0.96$ and $1.00$; actual change is $0.04$ and the bound $0.15$. An original absolute score strictly greater than $0.15$ is sufficient to preserve its sign.
\item Mean times are $14$ and $5.9$ milliseconds; the percentiles are $50$ and $5$. Compare each model's error conditional on its assigned route.
\end{enumerate}

\section{Challenges}

\begin{enumerate}
\item \textbf{Prove the allocation rule.} Under equal resource use and additive loss, use an exchange argument to show that the largest positive gains solve the fixed-budget problem. Explain why dependent labels do not invalidate the expectation calculation, while interference can invalidate the loss model.

\item \textbf{Derive complete selective risk.} Express complete loss using coverage, conditional model risk, conditional fallback risk, and referral cost. Construct two systems with the same coverage and model selective risk but different complete risk. Explain why a calibrated model score cannot resolve the difference by itself.

\item \textbf{Analyze distillation's limit.} Derive the softened-logit gradient and its high-temperature centered-logit limit. State the fixed-logit assumption and explain why this calculation does not imply an accuracy or latency guarantee.

\item \textbf{Separate observation from intervention.} A shadow model refers $40\%$ fewer cases. A team claims it would cut total waiting time by $40\%$. Identify what the shadow run establishes, what the queue model still needs, and how a shared queue could interfere with a randomized comparison.

\item \textbf{Compare resource-saving mechanisms.} Design an evaluation of a distilled model, a quantized model, and a cascade at a matched resource budget. Include outcomes on relevant slices, tail latency, routing errors, and the representation assumptions needed for each comparison.

\item \textbf{Reason about a changing policy source.} Institutional policy changes daily, some outcomes are delayed, and a retrieved page may contain hostile instructions. Define what an evaluated system version includes, which actions remain authorized, and how evidence updates, outcome audits, and fallback capacity interact. Identify a claim the available observations cannot yet support.
\end{enumerate}

\clearpage
{\pagestyle{empty}\cleardoublepage}
\endgroup
